% =============================================================
%  02_funzioni_dati.tex
%  Capitolo 2 --- Funzioni e dati fondamentali
% =============================================================

\chapter{Funzioni e dati fondamentali}

\begin{intuizione}
Prima di scrivere algoritmi complessi, devi padroneggiare
il vocabolario di base: i tipi di dato, i binding locali,
i condizionali e le varianti di \fn{defun}.
Questo capitolo è il lessico da cui parte tutto il resto.
\end{intuizione}

\vspace{4pt}

% ================================================================
\section{Tipi di dato}
% ================================================================

Common Lisp è a \textbf{tipizzazione dinamica}: ogni valore ha un tipo,
ma le variabili no. Il tipo viene controllato a runtime.

\argomento{Numeri}

\begin{concetto}{Numeri in Common Lisp}
Quattro sottotipi principali:
\begin{itemize}
  \item \textbf{Interi} (\texttt{integer}): dimensione arbitraria, nessun overflow
  \item \textbf{Razionali} (\texttt{ratio}): frazioni esatte, tipo \texttt{3/4}
  \item \textbf{Float} (\texttt{float}): virgola mobile IEEE 754
  \item \textbf{Complessi} (\texttt{complex}): raramente usati
\end{itemize}
\end{concetto}

\begin{esempio}[Aritmetica di base]
\begin{lstlisting}
(+ 3 4)          ; => 7      intero
(/ 10 3)         ; => 10/3   razionale esatto!
(/ 10.0 3)       ; => 3.3333... float
(* 1/3 3)        ; => 1      esatto

(expt 2 10)      ; => 1024
(sqrt 2)         ; => 1.4142...
(abs -7)         ; => 7
(mod 17 5)       ; => 2
(rem -7 3)       ; => -1     (rem mantiene segno del dividendo)

(floor  7/2)     ; => 3      (arrotonda verso -inf)
(ceiling 7/2)    ; => 4      (arrotonda verso +inf)
(truncate 7/2)   ; => 3      (taglia la parte decimale)
(round  7/2)     ; => 4      (arrotonda al piu' vicino, 0.5 -> pari)
\end{lstlisting}
\end{esempio}

\begin{attenzione}
\textbf{Divisione intera / non esiste come operatore separato.}
\texttt{(/ 10 3)} restituisce \texttt{10/3} (razionale), non \texttt{3}.
Per il quoziente intero usa \texttt{(floor (/ 10 3))} oppure
\texttt{(truncate 10 3)}, che con due argomenti restituisce quoziente
e resto come valori multipli.
\end{attenzione}

\seprule

\argomento{Simboli, booleani, caratteri, stringhe}

\begin{esempio}[Gli altri tipi fondamentali]
\begin{lstlisting}
;; Booleani
t              ; => T   (vero)
nil            ; => NIL (falso -- e lista vuota -- e simbolo NIL)
(not nil)      ; => T
(not 0)        ; => NIL  (0 e' VERO! Solo nil e' falso)
(not "")       ; => NIL  (stringa vuota e' VERA)
(not '())      ; => T    ('() e' nil, quindi e' FALSO)

;; Simboli
'hello         ; => HELLO  (case-insensitive)
'Hello         ; => HELLO  (identico)
(symbolp 'x)   ; => T

;; Caratteri
#\a            ; => #\a    (il carattere a)
#\Space        ; => #\Space
#\Newline      ; => #\Newline
(char-code #\A)   ; => 65
(code-char 65)    ; => #\A

;; Stringhe
"ciao"            ; => "ciao"
(length "ciao")   ; => 4
(string-upcase "ciao")   ; => "CIAO"
(string-downcase "CIAO") ; => "ciao"
(concatenate 'string "a" "b" "c")  ; => "abc"
\end{lstlisting}
\end{esempio}

\begin{attenzione}
\textbf{In Common Lisp, solo \texttt{nil} è falso. Tutto il resto è vero.}
Lo zero, la stringa vuota, un numero negativo --- tutto è vero.
La lista vuota \texttt{()} invece è falsa, perché \texttt{()} \emph{è} \texttt{nil}.
Questa è una differenza critica rispetto a Python o C.
\end{attenzione}

\begin{intuizione}
\texttt{nil} fa tre lavori contemporaneamente: è il valore falso,
è la lista vuota \texttt{()}, ed è il simbolo \texttt{NIL}.
Questi tre ruoli sono la stessa identità per design --- non tre cose
diverse, ma un unico oggetto con tre nomi.
\texttt{(null nil)} $\Rightarrow$ \texttt{T},
\texttt{(not nil)} $\Rightarrow$ \texttt{T},
\texttt{(listp nil)} $\Rightarrow$ \texttt{T},
\texttt{(equal nil '())} $\Rightarrow$ \texttt{T}.
\end{intuizione}

\seprule

\argomento{Predicati di tipo}

\begin{sintassi}
(numberp x)    (integerp x)   (floatp x)\\
(stringp x)    (symbolp x)    (characterp x)\\
(listp x)      (consp x)      (null x)\\
(functionp x)  (atom x)
\end{sintassi}

\begin{esempio}[Predicati di tipo]
\begin{lstlisting}
(numberp 42)       ; => T
(numberp "42")     ; => NIL
(atom 'x)          ; => T    (simbolo = atomo)
(atom '(a b))      ; => NIL  (lista = non atomo)
(consp '(a b))     ; => T    (lista non vuota)
(consp '())        ; => NIL  (nil non e' una cons-cell)
(null '())         ; => T
(null nil)         ; => T
(null 'x)          ; => NIL
\end{lstlisting}
\end{esempio}

% ================================================================
\section{Binding locale: \texttt{let} e \texttt{let*}}
% ================================================================

\begin{concetto}{\texttt{let} --- binding parallelo}
\begin{sintassi}
(let ((var1 val1)\\
\hspace{4em}(var2 val2)\\
\hspace{4em}...)\\
\hspace{2em}corpo)
\end{sintassi}
Tutte le espressioni \texttt{val} vengono valutate \textbf{prima} di
legare qualsiasi variabile. I binding sono quindi \emph{paralleli}:
\texttt{var1} non è ancora definita quando si valuta \texttt{val2}.
\end{concetto}

\begin{concetto}{\texttt{let*} --- binding sequenziale}
Identico a \texttt{let}, ma i binding avvengono \textbf{in sequenza}:
ogni variabile è visibile nelle espressioni successive.
Usalo quando un valore dipende da quello precedente.
\end{concetto}

\begin{esempio}[let vs let*]
\begin{lstlisting}
;; let: tutti i valori si calcolano nell'env ESTERNO
(let ((x 10)
      (y 20))
  (+ x y))         ; => 30

;; let*: y vede gia' x
(let* ((x 10)
       (y (* x 2))
       (z (+ x y)))
  z)               ; => 30

;; let NON funzionerebbe qui:
;; (let ((x 10) (y (* x 2))) ...)  ; ERRORE: x non ancora definita
\end{lstlisting}
\end{esempio}

\mnota{Regola pratica: usa \fn{let} se i binding sono indipendenti, \fn{let*} se uno dipende da un altro.}

\begin{esempio}[Pattern tipico: valori intermedi]
\begin{lstlisting}
(defun ipotenusa (a b)
  (let* ((a2 (* a a))
         (b2 (* b b))
         (somma (+ a2 b2)))
    (sqrt somma)))

(ipotenusa 3 4)   ; => 5.0
\end{lstlisting}
\end{esempio}

\begin{attenzione}
\textbf{Le variabili di \texttt{let} non esistono fuori dal corpo.}
Sono locali all'espressione \texttt{let}. Non "escono".
Non è necessario --- e non si deve --- dichiararle prima.
Questo è radicalmente diverso da C o Java.
\end{attenzione}

\begin{workbook}
\textbf{Esercizi 4--5} (Ipotenusa, Distanza euclidea)\\[4pt]
Entrambi richiedono \fn{let} o \fn{let*} per i valori intermedi.
L'esercizio 4 impone esplicitamente \fn{let}: è un esercizio
\emph{didattico} sul binding, non solo sull'aritmetica.
\end{workbook}

% ================================================================
\section{Condizionali}
% ================================================================

\argomento{\texttt{if} --- la forma base}

\begin{concetto}{\texttt{if}}
\begin{sintassi}
(if condizione\\
\hspace{2em}espressione-se-vera\\
\hspace{2em}espressione-se-falsa)  ; opzionale
\end{sintassi}
\texttt{if} è un'\textbf{espressione}: restituisce il valore del ramo
eseguito. Se il ramo falso manca e la condizione è falsa,
restituisce \texttt{NIL}.
\end{concetto}

\begin{esempio}[\texttt{if} come espressione]
\begin{lstlisting}
(if (> 3 2) "maggiore" "minore")   ; => "maggiore"

;; Usabile direttamente in un'espressione piu' grande
(+ 10 (if (oddp 7) 1 0))   ; => 11

;; if con una sola branch
(if (> x 0) (print "positivo"))
;; => NIL se x <= 0, stampa e restituisce "positivo" altrimenti
\end{lstlisting}
\end{esempio}

\begin{nota}
\fn{if} accetta \textbf{un'unica espressione} per branch.
Se vuoi eseguire più espressioni in un branch, usa \fn{progn}:
\texttt{(if cond (progn e1 e2 e3) altro)}.
In pratica, quando il corpo è lungo si preferisce \fn{when}/\fn{unless}
o \fn{cond}.
\end{nota}

\seprule

\argomento{\texttt{when} e \texttt{unless} --- versioni semplificate}

\begin{sintassi}
(when condizione corpo...)    ; esegue corpo se vero, ritorna NIL altrimenti\\
(unless condizione corpo...)  ; esegue corpo se FALSO
\end{sintassi}

\begin{esempio}[\texttt{when}/\texttt{unless}]
\begin{lstlisting}
(when (> x 0)
  (print "positivo")
  (print "elaboro..."))   ; piu' espressioni ok

(unless (null lst)
  (print (car lst)))
\end{lstlisting}
\end{esempio}

\seprule

\argomento{\texttt{cond} --- ramificazione multipla}

\begin{concetto}{\texttt{cond}}
\begin{sintassi}
(cond (test1 expr1...)\\
\hspace{2em}  (test2 expr2...)\\
\hspace{2em}  ...\\
\hspace{2em}  (t     default...))
\end{sintassi}
Valuta i test in ordine. Appena uno è vero, esegue le espressioni
associate e restituisce l'ultima. La clausola \texttt{(t ...)} è
il caso default (opzionale ma raccomandato).
\end{concetto}

\begin{esempio}[\texttt{cond} per ramificazioni multiple]
\begin{lstlisting}
(defun segno (n)
  (cond ((< n 0) 'negativo)
        ((= n 0) 'zero)
        (t       'positivo)))

(segno -5)   ; => NEGATIVO
(segno 0)    ; => ZERO
(segno 7)    ; => POSITIVO


;; Voto in lettere
(defun voto->giudizio (p)
  (cond ((>= p 90) "ottimo")
        ((>= p 75) "buono")
        ((>= p 60) "sufficiente")
        (t         "insufficiente")))
\end{lstlisting}
\end{esempio}

\begin{intuizione}
\textbf{Quando usare cosa:}
\begin{itemize}
  \item \fn{if} --- due casi, espressioni brevi
  \item \fn{when}/\fn{unless} --- un solo caso, corpo multi-espressione
  \item \fn{cond} --- tre o più casi, o condizioni complesse
\end{itemize}
\end{intuizione}

\seprule

\argomento{\texttt{case} --- dispatch su valori discreti}

\begin{sintassi}
(case valore\\
\hspace{2em}(chiave1 expr...)\\
\hspace{2em}(chiave2 expr...)\\
\hspace{2em}(otherwise default...))
\end{sintassi}

\begin{esempio}[\texttt{case}]
\begin{lstlisting}
(defun calcola (op a b)
  (case op
    (+  (+ a b))
    (-  (- a b))
    (*  (* a b))
    (/  (if (= b 0) nil (/ a b)))
    (otherwise (error "operatore sconosciuto"))))

(calcola '+ 3 4)   ; => 7
(calcola '* 5 6)   ; => 30
\end{lstlisting}
\end{esempio}

\begin{nota}
\fn{case} confronta con \fn{eql} (uguaglianza per simboli e numeri).
Non funziona con stringhe --- per quelle usa \fn{cond} con \fn{equal}.
\end{nota}

\begin{workbook}
\textbf{Esercizi 8--16} (Segno, Massimo, Triangolo, Bisestile,
Giorni mese, Voto, Terna, Radici, Calcolatrice)\\[4pt]
Questi esercizi coprono \fn{cond}, \fn{case}, \fn{and}/\fn{or} e \fn{let*}.
L'ordine suggerito: 8, 11, 15, 16 per gli essenziali;
9, 10, 12, 13, 14 per completare.\\[4pt]
L'esercizio 15 (radici quadratiche) richiede \fn{let*} e la gestione
del caso $\Delta < 0$: è il primo esercizio che combina binding e
condizionali in modo non banale.
\end{workbook}

% ================================================================
\section{Operatori logici: \texttt{and}, \texttt{or}, \texttt{not}}
% ================================================================

\begin{concetto}{\texttt{and}, \texttt{or} --- short-circuit e valori}
\texttt{and} e \texttt{or} non sono solo operatori booleani:
restituiscono \emph{valori}, non solo \texttt{T}/\texttt{NIL}.

\medskip
\begin{tabular}{@{}ll@{}}
\texttt{(and e1 e2 ...)} & Restituisce il primo falso, o l'ultimo valore\\
\texttt{(or e1 e2 ...)}  & Restituisce il primo vero, o \texttt{NIL}\\
\texttt{(not x)}         & Restituisce \texttt{T} o \texttt{NIL}\\
\end{tabular}
\medskip

Entrambi usano la \textbf{valutazione cortocircuito}: si fermano
appena conoscono il risultato.
\end{concetto}

\begin{esempio}[\texttt{and}/\texttt{or} come flusso di controllo]
\begin{lstlisting}
(and 1 2 3)          ; => 3      (tutti veri: restituisce l'ultimo)
(and 1 nil 3)        ; => NIL    (si ferma al primo falso)
(or  nil nil 42)     ; => 42     (primo vero)
(or  nil nil nil)    ; => NIL
(or  nil "default")  ; => "default"

;; Uso idiomatico: valore di default
(defun saluta (nome)
  (let ((n (or nome "Utente")))
    (format t "Ciao, ~a!~%" n)))

(saluta nil)    ; stampa: Ciao, Utente!
(saluta "Ada")  ; stampa: Ciao, Ada!

;; Uso idiomatico: guard clause
(defun dividi (a b)
  (and (not (= b 0))
       (/ a b)))

(dividi 10 2)   ; => 5
(dividi 10 0)   ; => NIL
\end{lstlisting}
\end{esempio}

% ================================================================
\section{Parametri opzionali, con default e multipli}
% ================================================================

\argomento{\texttt{\&optional} --- parametri facoltativi}

\begin{concetto}{\texttt{\&optional}}
\begin{sintassi}
(defun f (obbligatorio \&optional (opt1 default1) (opt2 default2))\\
\hspace{2em}corpo)
\end{sintassi}
Se non si passa il parametro, viene usato il default.
Se il default è omesso, il valore è \texttt{NIL}.
\end{concetto}

\begin{esempio}[\texttt{\&optional}]
\begin{lstlisting}
(defun saluta (&optional (nome "Mondo"))
  (format t "Ciao, ~a!~%" nome))

(saluta)        ; stampa: Ciao, Mondo!
(saluta "Ada")  ; stampa: Ciao, Ada!


(defun potenza (base &optional (esp 2))
  (expt base esp))

(potenza 3)      ; => 9   (3^2)
(potenza 3 4)    ; => 81  (3^4)
\end{lstlisting}
\end{esempio}

\seprule

\argomento{\texttt{\&rest} --- numero variabile di argomenti}

\begin{concetto}{\texttt{\&rest}}
\begin{sintassi}
(defun f (a b \&rest resto) corpo)
\end{sintassi}
\texttt{resto} è una \emph{lista} di tutti gli argomenti extra.
\end{concetto}

\begin{esempio}[\texttt{\&rest}]
\begin{lstlisting}
(defun somma-tutto (&rest numeri)
  (apply #'+ numeri))

(somma-tutto 1 2 3 4 5)   ; => 15
(somma-tutto)              ; => 0


(defun mio-format (fmt &rest args)
  (apply #'format t fmt args))
\end{lstlisting}
\end{esempio}

\seprule

\argomento{Valori multipli: \texttt{values} e \texttt{multiple-value-bind}}

\begin{concetto}{Valori multipli}
Una funzione Lisp può restituire \textbf{più valori} contemporaneamente
--- senza costruire una lista.
\end{concetto}

\begin{esempio}[Valori multipli]
\begin{lstlisting}
;; Restituire due valori
(defun min-max (lst)
  (values (apply #'min lst)
          (apply #'max lst)))

;; Ricevere piu' valori
(multiple-value-bind (minimo massimo)
    (min-max '(3 1 4 1 5 9 2 6))
  (format t "min=~a max=~a~%" minimo massimo))
; stampa: min=1 max=9


;; floor con 2 argomenti restituisce quoziente e resto
(multiple-value-bind (q r)
    (floor 17 5)
  (format t "~a = ~a * 5 + ~a~%" 17 q r))
; stampa: 17 = 3 * 5 + 2
\end{lstlisting}
\end{esempio}

\begin{nota}
Se usi una funzione multi-valore in un contesto che ne vuole uno solo
(es. \texttt{(+ (floor 7 2) 1)}), viene usato solo il primo valore.
\end{nota}

% ================================================================
\section{Iterazione di base}
% ================================================================

La ricorsione è il modo idiomatico in Lisp. Ma Lisp ha anche
costrutti iterativi espliciti, utili per output, effetti collaterali
e loop semplici.

\argomento{\texttt{dotimes} --- ripetizione N volte}

\begin{concetto}{\texttt{dotimes}}
\begin{sintassi}
(dotimes (variabile N risultato)\\
\hspace{2em}corpo)
\end{sintassi}
\texttt{variabile} va da \texttt{0} a \texttt{N-1}.
\texttt{risultato} è il valore restituito (default \texttt{NIL}).
\end{concetto}

\begin{esempio}[\texttt{dotimes}]
\begin{lstlisting}
(dotimes (i 5)
  (format t "~a " i))
; stampa: 0 1 2 3 4


;; Accumulatore con dotimes
(let ((somma 0))
  (dotimes (i 10 somma)           ; restituisce somma
    (setf somma (+ somma i))))
; => 45
\end{lstlisting}
\end{esempio}

\seprule

\argomento{\texttt{dolist} --- iterare su una lista}

\begin{sintassi}
(dolist (variabile lista risultato)\\
\hspace{2em}corpo)
\end{sintassi}

\begin{esempio}[\texttt{dolist}]
\begin{lstlisting}
(dolist (x '(1 2 3 4 5))
  (format t "~a " (* x x)))
; stampa: 1 4 9 16 25
\end{lstlisting}
\end{esempio}

\seprule

\argomento{\texttt{loop} --- iterazione flessibile}

\fn{loop} è il costrutto più potente. Nella forma base:

\begin{esempio}[\texttt{loop} base]
\begin{lstlisting}
;; Iterazione con contatore
(loop for i from 1 to 5
      do (format t "~a " i))
; stampa: 1 2 3 4 5


;; Raccogliere risultati
(loop for i from 1 to 5
      collect (* i i))
; => (1 4 9 16 25)


;; Con accumulatore
(loop for i from 1 to 100
      sum i)
; => 5050


;; Iterare su una lista
(loop for x in '(a b c d)
      do (print x))
\end{lstlisting}
\end{esempio}

\begin{attenzione}
\textbf{\fn{setf} e mutabilita'.}
Negli esempi sopra compare \fn{setf}: è l'operatore di assegnazione.
Usarlo è lecito, ma nella programmazione funzionale si preferisce
la ricorsione con accumulatori. Ritorneremo su questo nel
Capitolo~4 (Higher-order e closure).

In generale: \fn{dotimes}/\fn{dolist}/\fn{loop} per output e side-effect;
ricorsione per trasformare dati.
\end{attenzione}

\begin{workbook}
\textbf{Esercizi 17--21} (Somma cifre, Conta fino a N, Tabellina,
Sconto, Paga settimanale)\\[4pt]
L'esercizio 17 (Somma delle cifre) è il più interessante:
richiede ricorsione su un numero intero, estraendo le cifre
con \fn{mod} e \fn{floor}. Ha un suggerimento nel workbook.\\[4pt]
Gli esercizi 18 e 19 sono esplicitamente su \fn{dotimes}/\fn{loop}:
l'obiettivo è prendere confidenza con l'output iterativo.
\end{workbook}

% ================================================================
\section{Il formato dell'output: \texttt{format}}
% ================================================================

\fn{format} è la funzione di output principale. La forma base:

\begin{sintassi}
(format destinazione stringa-controllo argomenti...)
\end{sintassi}

\texttt{t} = standard output; \texttt{nil} = restituisce una stringa.

\begin{esempio}[Direttive \texttt{format} essenziali]
\begin{lstlisting}
(format t "Ciao, ~a!~%"  "Ada")   ; ~a = qualsiasi valore
(format t "Valore: ~s~%" '(1 2))  ; ~s = rappresentazione leggibile
(format t "Numero: ~d~%" 42)      ; ~d = intero decimale
(format t "Float: ~f~%"  3.14)    ; ~f = float
(format t "~%")                   ; ~% = newline
(format t "~&")                   ; ~& = newline se non gia' a inizio riga
(format t "~10d~%" 42)            ; ~10d = campo di 10 caratteri

;; Restituire stringa invece di stampare
(format nil "Risultato: ~a" 42)   ; => "Risultato: 42"
\end{lstlisting}
\end{esempio}

% ================================================================
\section*{Riepilogo del capitolo}
% ================================================================

\begin{tcolorbox}[
  enhanced,
  colback=csBlue!5, colframe=csBlue!40,
  arc=4pt, boxrule=0.6pt,
  left=14pt, right=14pt, top=12pt, bottom=12pt,
  before skip=16pt, after skip=16pt
]
\textbf{Quello che hai imparato:}

\medskip
\begin{itemize}
  \item \textbf{Tipi}: interi senza overflow, razionali esatti,
        float, simboli, stringhe, caratteri.
        Solo \texttt{NIL} è falso.
  \item \textbf{\texttt{let}/\texttt{let*}}: binding locale.
        Parallelo vs sequenziale.
        Le variabili non "escono" dal corpo.
  \item \textbf{Condizionali}: \texttt{if} (2 casi),
        \texttt{when}/\texttt{unless} (1 caso + corpo multiplo),
        \texttt{cond} (n casi), \texttt{case} (dispatch su valori).
  \item \textbf{\texttt{and}/\texttt{or}}: short-circuit con valori reali,
        non solo booleani.
  \item \textbf{Parametri}: \texttt{\&optional} (con default),
        \texttt{\&rest} (argomenti variabili).
  \item \textbf{Valori multipli}: \texttt{values} +
        \texttt{multiple-value-bind}.
  \item \textbf{Iterazione}: \texttt{dotimes}, \texttt{dolist},
        \texttt{loop} --- per output e side-effect.
  \item \textbf{\texttt{format}}: output formattato con direttive
        \texttt{\textasciitilde a}, \texttt{\textasciitilde d},
        \texttt{\textasciitilde f}, \texttt{\textasciitilde \%}.
\end{itemize}

\medskip
\textbf{Cosa viene dopo:}\\
Capitolo~3 --- \emph{Condizioni e flusso avanzato}: ricorsione
strutturata, tail recursion, gestione degli errori, \texttt{block}/\texttt{return-from}.
\end{tcolorbox}

\begin{workbook}
\textbf{Prima di andare al Capitolo~3:}\\[4pt]
Completa gli esercizi \textbf{1--21} del workbook (capitolo
\emph{Fondamenti}).
Sono tutti fattibili con i costrutti di questo capitolo.\\[4pt]
Se sei a corto di tempo: priorita' agli esercizi essenziali
(contrassegnati con $\blacklozenge$): 1, 4, 6, 8, 11, 15, 17, 18.
\end{workbook}
